\documentclass[tikz,border=4pt]{standalone}
\usepackage{ctex}
\usepackage{amsmath}
\usetikzlibrary{shapes.geometric, arrows.meta, positioning}

\begin{document}
\begin{tikzpicture}[
  node distance=1.1cm,
  box/.style={rectangle, draw=black, rounded corners=3pt,
              minimum width=3.2cm, minimum height=0.65cm,
              align=center, font=\small},
  dbox/.style={diamond, draw=black, aspect=2.5,
               minimum width=3.5cm, minimum height=1cm,
               align=center, font=\small},
  arr/.style={->, >=Stealth, thick},
]

  \node[box, fill=blue!10, minimum width=6cm] (start)
    {含参方程组 $\begin{cases}a_1 x+b_1 y=c_1\\a_2 x+b_2 y=c_2\end{cases}$};

  \node[box, below=of start] (step1)
    {比较系数比 $\dfrac{a_1}{a_2}$ 与 $\dfrac{b_1}{b_2}$};

  \node[dbox, below=1.2cm of step1, fill=yellow!20] (cond1)
    {$\dfrac{a_1}{a_2} \neq \dfrac{b_1}{b_2}$？};

  \node[box, below left=1.2cm and 1.5cm of cond1, fill=green!15] (unique)
    {两直线相交\\唯一解};

  \node[dbox, below right=1.2cm and 1.2cm of cond1, fill=yellow!20] (cond2)
    {$\dfrac{c_1}{c_2} = \dfrac{a_1}{a_2}$？};

  \node[box, below left=1.2cm and 0.5cm of cond2, fill=red!15] (nosol)
    {两直线平行\\无解};

  \node[box, below right=1.2cm and 0.5cm of cond2, fill=cyan!15] (inf)
    {两直线重合\\无穷多解};

  \draw[arr] (start) -- (step1);
  \draw[arr] (step1) -- (cond1);
  \draw[arr] (cond1) -- node[above left, font=\small]{是} (unique);
  \draw[arr] (cond1) -- node[above right, font=\small]{否} (cond2);
  \draw[arr] (cond2) -- node[above left, font=\small]{否} (nosol);
  \draw[arr] (cond2) -- node[above right, font=\small]{是} (inf);

\end{tikzpicture}
\end{document}
